\apendice{Anexo de sostenibilización curricular}

\section{Introducción}

La sostenibilidad del proyecto no depende únicamente de que la aplicación se ejecute en la nube. También intervienen el número de recursos desplegados, su tamaño, la posibilidad de mantener el software y la utilidad del proceso automatizado \cite{crue_sostenibilidad_directrices,azure_cloud_adoption_framework}. Por esa razón, esta reflexión distingue entre efectos observables del prototipo y beneficios que no se han medido. Los ODS más relacionados son el ODS 9, industria, innovación e infraestructura; el ODS 12, producción y consumo responsables; y, de forma indirecta, el ODS 13, acción por el clima.

Respecto al ODS 9, el resultado es una infraestructura digital que puede revisarse y reproducirse. App Service ejecuta la aplicación, Blob Storage conserva las solicitudes y Key Vault separa los secretos del código. Terraform y Bicep documentan la forma de los recursos, mientras que los scripts PowerShell y Azure CLI recogen el despliegue que se ejecutó realmente. Esta combinación evita que el funcionamiento dependa de una secuencia de acciones recordada solo por la autora y facilita auditar la configuración \cite{azure_well_architected}.

En relación con el ODS 12, se evitó reproducir una arquitectura corporativa sobredimensionada. El prototipo utiliza una sola región, un plan App Service B1, almacenamiento con replicación local y un único entorno de desarrollo. Tampoco se desplegaron máquinas virtuales, clústeres ni bases de datos que el caso práctico no necesitaba. Estas decisiones reducen coste y número de componentes, aunque no convierten por sí mismas la solución en sostenible. Los recursos deberán eliminarse cuando termine su periodo de evaluación para evitar consumo y cargos sin utilidad.

La relación con el ODS 13 es más limitada. El proyecto no dispone de medidas de consumo eléctrico ni de emisiones y, por tanto, no puede afirmar una reducción cuantificada frente a una instalación local. Sí puede señalar que usa servicios compartidos y evita adquirir hardware específico para la prueba. La contrapartida es que App Service, Storage, Key Vault y Monitor permanecen en centros de datos que consumen energía. Una continuación rigurosa debería registrar tiempos de uso, región, tamaño de los recursos y estimaciones de carbono antes de comparar alternativas.

La dimensión organizativa se observa con más claridad. Cada solicitud recibe un identificador, estado, prioridad y recomendación. El equipo de TI puede consultar el conjunto mediante un endpoint protegido y revisar agregados básicos. Frente a una cadena de correos, este registro reduce la posibilidad de perder una petición y hace visible el criterio usado para priorizarla. El prototipo no mide tiempos de resolución, por lo que una mejora de esos tiempos sigue siendo una hipótesis que requeriría datos de uso reales.

La seguridad forma parte de la sostenibilidad porque un sistema difícil de proteger tampoco resulta mantenible. Key Vault guarda la clave de consulta y Managed Identity permite que App Service la recupere sin otra credencial intermedia \cite{azure_security_design}. Los ejemplos utilizan identidades ficticias y no se ha cargado información personal real. Aun así, una puesta en producción necesitaría identidad individual, autorización por roles, política de conservación y un análisis formal de protección de datos.

La mantenibilidad se apoya en siete pruebas automáticas, documentación de endpoints, diagramas, decisiones de arquitectura y scripts de verificación. Otra persona puede localizar los componentes y repetir las comprobaciones principales sin depender de una explicación oral. Esto no garantiza por sí solo la continuidad del software, pero reduce el esfuerzo necesario para entender qué se desplegó y por qué.

En conjunto, la aportación sostenible más defendible es la moderación del diseño: desplegar solo los servicios que apoyan un requisito, documentar cómo retirarlos o reproducirlos y reconocer los impactos que no se han medido. El prototipo ofrece una base mantenible y segura para estudiar el proceso, pero una valoración ambiental completa exigiría métricas que quedan fuera del alcance de este TFG.
